% latex-ops example: Problem Set + Swiss Minimal, solutions version.
% Prompt: "Problem set on eigenvalues for a first-year linear algebra course, 4 problems, with
% solutions in a separate version from the same source."
% Every solution below was checked by hand and numerically (characteristic polynomials and eigenvectors).
\documentclass[11pt,a4paper]{article}

\usepackage[english]{babel}
\usepackage{fontspec}
\usepackage{amsmath,mathtools}
\usepackage{unicode-math}
\usepackage{microtype}
\usepackage[a4paper,margin=24mm]{geometry}
\usepackage{xcolor}
\usepackage{enumitem}
\usepackage[hidelinks]{hyperref}

\IfFontExistsTF{Inter}{\setsansfont{Inter}}{\setsansfont{Latin Modern Sans}}
\setmainfont{Latin Modern Roman}
\setmathfont{Latin Modern Math}
\definecolor{Accent}{HTML}{315B6B}
\definecolor{Muted}{HTML}{6B6B6B}
\setlength{\parindent}{0pt}
\setlength{\parskip}{.45em}
\setlist{parsep=0pt,itemsep=.3em}
\widowpenalty=10000 \clubpenalty=10000
\pagestyle{plain}

\newif\ifsolutions \solutionstrue % \solutionsfalse prints answer space instead
\newcounter{problem}
\newenvironment{problem}[1]
  {\refstepcounter{problem}\par\bigskip\noindent\begin{minipage}{\linewidth}\setlength{\parskip}{.45em}%
   {\sffamily\bfseries Problem \theproblem}\hfill{\sffamily\small\color{Muted}#1 points}\par\smallskip}
  {\end{minipage}\par}
\newcommand{\answerspace}[1]{\ifsolutions\else\par\vspace{#1}\fi}
\newenvironment{solution}
  {\ifsolutions\par\medskip\noindent{\sffamily\bfseries\color{Accent}Solution.}\ \else\setbox0\vbox\bgroup\fi}
  {\ifsolutions\par\else\egroup\fi}

\begin{document}

{\sffamily\small\color{Muted}LINEAR ALGEBRA I \textperiodcentered\ WEEK 9\par}
{\sffamily\LARGE\bfseries Problem Set 5: Eigenvalues\ifsolutions{} --- Solutions\fi\par}
{\sffamily\color{Muted}Due Friday \textperiodcentered\ 20 points\par}

\medskip
\textbf{Instructions.} Show every step. A correct answer without justification earns no credit.

\begin{problem}{5}
Find the eigenvalues and an eigenvector for each eigenvalue of
$A = \begin{pmatrix} 2 & 1 \\ 1 & 2 \end{pmatrix}$.
\answerspace{5cm}
\begin{solution}
$\det(A - \lambda I) = (2-\lambda)^2 - 1 = (\lambda - 1)(\lambda - 3)$, so $\lambda_1 = 1$ and
$\lambda_2 = 3$. For $\lambda_1 = 1$, $(A - I)v = 0$ gives $v_1 = (1, -1)^\top$; for
$\lambda_2 = 3$, $(A - 3I)v = 0$ gives $v_2 = (1, 1)^\top$.
\end{solution}
\end{problem}

\begin{problem}{5}
Show that if $\lambda$ is an eigenvalue of an invertible matrix $B$, then $\lambda \neq 0$ and
$1/\lambda$ is an eigenvalue of $B^{-1}$.
\answerspace{5cm}
\begin{solution}
If $\lambda = 0$ then $Bv = 0$ for some $v \neq 0$, contradicting invertibility. From $Bv = \lambda v$,
multiply by $B^{-1}$: $v = \lambda B^{-1} v$, so $B^{-1} v = \tfrac{1}{\lambda} v$ with the same $v \neq 0$.
\end{solution}
\end{problem}

\begin{problem}{5}
Let $C = \begin{pmatrix} 0 & -1 \\ 1 & 0 \end{pmatrix}$. Show that $C$ has no real eigenvalues and
find its complex eigenvalues.
\answerspace{4cm}
\begin{solution}
$\det(C - \lambda I) = \lambda^2 + 1$, which has no real root. Over $\mathbb{C}$ the eigenvalues
are $\lambda = \pm i$. Geometrically, $C$ rotates the plane by $90^\circ$, so no real direction is preserved.
\end{solution}
\end{problem}

\begin{problem}{5}
The trace of a $2 \times 2$ matrix $M$ is $5$ and its determinant is $6$. Find the eigenvalues of $M$.
\answerspace{4cm}
\begin{solution}
The characteristic polynomial is $\lambda^2 - (\operatorname{tr} M)\lambda + \det M = \lambda^2 - 5\lambda + 6
= (\lambda - 2)(\lambda - 3)$, so the eigenvalues are $2$ and $3$.
\end{solution}
\end{problem}

\vfill
{\small\color{Muted}Illustrative example generated with latex-ops.\par}
\end{document}
